\scanpage{59}
\unsure{10|18}. Cas mixte.

\uline{Th}. \deleted{\unsure{Soient}} Soit $S$ un schéma tel que $2\cdot1_S$ soit $\underline O_S$-régulier, \deleted{soit $S_0$} \inserted{\deleted{soit} $S_0=\deleted{\unsure{S}}V(2\underline O_S)$ $S_1=V(4\cdot\underline O_S)$.}

Pour tout rev. quadratique $\Delta$ de $S$, considérons le triple $(L,D,b_0)$ formé du module inversible $L\simeq\underline{\det}(\underline{\mathrm{Aff}}(\Delta/S))^{-1}$, de $\deleted{\unsure{\Theta|\partial}}=D(\Delta/S)\in\Gamma(S,L^{\otimes2})$, et de la section $b_0=b(\Delta_0/S_0)\in\Gamma(S_0,L_0)$, [où $\Delta_0=\Delta_{S_0},L=L_{S_0}$], de sorte que l'on a
\[b_0^2=D_0\tag*{(8.1)}\]

a) Le foncteur (\deleted{de} \inserted{entre} groupoïdes)
\[\Delta\longmapsto(L,D,b_0)\]
est pleinement fidèle.

b) \deleted{Supposons} Pour que \deleted{\unsure{tel}} $(L,D,b_0)$ appartienne à l'image essentielle du foncteur précédent, il \deleted{\unsure{suffit} et il faut} \inserted{faut, et il suffit} si $H^1(S,L)\longrightarrow H^1(S\setminus S_0,L)$ est injectif \inserted{(p.ex. si $S$ est affine\ldots)} \deleted{que $D_1=D|S_1$ soit de la forme $b_1^2$, où $b_1\in\Gamma(S_1,L_1)$ $(L_1=L_{S_1})$, avec $b_1|S_0=b_0$, [\unsure{i.e. $D_1=?$ $b_0$}] ou encore} \inserted{qu'on ait $D=b^2-4c$ avec $b\in\Gamma(S,L)$, $c\in\Gamma(S,L^{\otimes2})$.}

c) En tous les cas, pour que $(L,D,b_0)$ appartienne à l'image essentielle, il f et s que l'on ait
\[D_1=\gamma(b_0),\]
où
\[\gamma:\deleted{\unsure{L|S_0}}\ \mathcal O_{S_0}\longrightarrow\underline O_{S_1}\]
est le morphisme de faisceaux \inserted{multiplicatifs} déduit de l'endom.
\[\underline O_{S_1}\longrightarrow\underline O_{S_1},\quad\beta_1\longmapsto\beta_1^2.\]
[compte tenu que $\beta_1\equiv\beta'_1\ (2)\implies\beta_1^2\equiv{\beta'_1}^2\ (4)$].

\deleted{\uline{Dém} a) \unsure{OPS|DPS} \inserted{\unsure{sur|de}} $S$ affine. On note que \inserted{le torseur} $F$ \inserted{sous $L$}, en tant que faisceau \inserted{le prétors\unsure{eur}} Zariski, \deleted{\illegible{1}} \inserted{$F\subset\frac12 L$} est déterminé par $b_0$, comme le sous-faisceau \inserted{de} $\frac12 \deleted{\unsure{L}}\supset\underline L$ image}
% archive sheet number: 58
